\chapter{Marco Teórico y Trabajos Relacionados}
\label{ch:marco}

Este capítulo ubica la tesis dentro de cuatro cuerpos de literatura: la tradición
Nodo--Lugar, que concibe las estaciones de transporte como puntos donde coinciden una red
de transporte y un lugar urbano; la estimación de la demanda de viajes, que aporta la
señal de demanda sobre la que se construye el marco; el análisis de accesibilidad y de
brecha de cobertura, que convierte demanda y oferta en un diagnóstico; y la literatura de
aprendizaje automático y de diseño de redes de la que provienen los componentes
predictivo y generativo. El capítulo cierra situando el vacío de investigación ---una
metodología basada en la demanda, no circular y transferible--- que aborda el resto de la
tesis.

\section{El modelo Nodo--Lugar y sus extensiones}
\label{sec:marco-nodolugar}

El ancla conceptual de este trabajo es el modelo Nodo--Lugar de Bertolini
\parencite{bertolini1996, bertolini1999}, que caracteriza el área de una estación según
dos ejes: su valor de \emph{nodo}, la posición que ocupa en la red de transporte, y su
valor de \emph{lugar}, la intensidad y diversidad de las actividades urbanas a su
alrededor. La tesis central del modelo es que el desarrollo orientado al transporte
requiere que ambos se mantengan en equilibrio; un desequilibrio persistente señala una
estación subutilizada en un entorno urbano rico, o un nodo sobredimensionado en uno débil,
lectura que después se operacionaliza como una clasificación empírica de estaciones
\parencite{reusser2008, vale2015}. El propio eje de lugar se apoya en una amplia literatura
del entorno construido que relaciona densidad, diversidad de usos de suelo y diseño con el
comportamiento de viaje \parencite{cervero1997threeD, ewing2010meta}, lo que motiva los
indicadores de lugar usados aquí. Como diagnóstico resulta atractivo porque es descriptivo
y de bajos requerimientos de datos ---los índices de nodo y lugar se arman solo con
variables de red y de uso de suelo---, pero en su forma original es estático y nada dice
sobre la \emph{demanda latente} ni sobre dónde deberían colocarse nuevos nodos.

Extensiones recientes agregan dimensiones a las dos originales, generalmente acoplando el
conjunto ampliado de indicadores con aprendizaje automático. \textcite{su2022nodeplace}
añaden a nodo y lugar un término de afluencia y lo interpretan con \textsc{shap}, moviendo
la tradición de la descripción hacia la predicción; \textcite{aminipishro2022nprt} en
cambio añaden tanto afluencia como hora del día como cuarta y tercera dimensión (una
formulación NP--RT, no NP--RV), clasificando estaciones con un método K-means/Cubo.
Ninguno añade un término de \emph{vitalidad} específicamente, y no pudo localizarse
ningún artículo verificado que corresponda a una formulación ``NP--RV''
nodo--lugar--afluencia--vitalidad durante una revisión de literatura del 2026-09-09; una
versión anterior de este capítulo citaba tal artículo, el cual se retira aquí en lugar de
dejarse sin verificar. El patrón más amplio que ilustran los dos artículos verificados
anteriores ---aumentar Nodo--Lugar con una variable de resultado y un modelo
interpretable--- es el que sigue esta tesis estructuralmente, aunque se aparta de él en la
variable de resultado misma (\cref{sec:marco-aprendizaje}).

La propia extensión posterior de Bertolini empareja el modelo con una medida explícita de
accesibilidad a través de varias zonas metropolitanas europeas en lugar de un término de
afluencia/vitalidad \parencite{papa2015tod}, la cual es el antecedente de linaje más
cercano a la extensión basada en accesibilidad y brecha de cobertura desarrollada en esta
tesis (\cref{sec:w3}). Esta tesis adopta una variante
Nodo--Lugar--Personas (\nppv): conserva nodo y lugar, pone en primer plano una dimensión
de \emph{personas} (población, dependencia y atributos de equidad) y ---de manera
crucial--- retira el término de vitalidad de la ponderación de priorización porque, en el
área de estudio, el único proxy de vitalidad disponible está resuelto a nivel municipal y
por tanto colapsa a una bandera binaria de ``tiene transporte prémium'' que reintroduce
justo la circularidad de seguir a la oferta que la tesis busca evitar (\cref{sec:w4}). Esta
exclusión se reporta aquí como una limitación de disponibilidad de datos del proxy de
vitalidad específico disponible en \zmg, no como una afirmación teórica de que la vitalidad
sea prescindible en principio para un modelo de tipo Nodo--Lugar; un proxy de vitalidad a
resolución de \ageb (p. ej., a partir de conteos peatonales o una medida de intensidad de
actividad independiente de la demanda) sigue siendo una dirección abierta, y la etiqueta
\nppv{} renombrada debe leerse como acotada a esta aplicación y no como una afirmación
general de que una dimensión de afluencia o vitalidad sea prescindible para las
extensiones Nodo--Lugar en general.

\section{Estimación de la demanda de viajes}
\label{sec:marco-demanda}

La capa de demanda sigue la tradición clásica de cuatro pasos
\parencite{mcnally2000fourstep, ortuzar2011}, de la cual se usan aquí la generación y la
distribución de viajes. La distribución se modela con un modelo de interacción espacial
(gravitacional), cuya justificación estadística moderna es la derivación por maximización
de entropía de \textcite{wilson1971gravity}: entre todas las matrices origen--destino
consistentes con los totales marginales observados, la forma gravitacional es la más
probable, y su parámetro de decaimiento con la distancia gobierna cuán abruptamente cae la
interacción con la separación. La variante doblemente restringida, balanceada por el
procedimiento de ajuste proporcional iterativo (Furness) \parencite{furness1965}, reproduce
exactamente los totales de origen y destino y es la forma adoptada en \cref{sec:w1}.

Dos decisiones de diseño distinguen esta capa de demanda de una aplicación de libro de
texto. Primero, la \emph{generación} de viajes se construye por completo con variables
censales y económicas de Nivel~1, sin ningún insumo de transporte, de modo que la demanda
nunca es función de la oferta actual ---la propiedad de la que depende todo el marco.
Segundo, el parámetro de decaimiento no se supone sino que se \emph{calibra} contra una
encuesta metropolitana independiente \parencite{imeplan_eod2022}, para que la función de
impedancia refleje el comportamiento observado y no un exponente convencional
(\cref{sec:w2}). La calibración también arroja un error de transferencia documentado, que
importa cuando el mismo previo se traslada a otras ciudades (\cref{ch:transferibilidad}).

\section{Accesibilidad y análisis de brecha de cobertura}
\label{sec:marco-accesibilidad}

La oferta se mide con una superficie de accesibilidad de oportunidades acumuladas: para
cada unidad, el número de oportunidades (aquí, empleos) alcanzables dentro de un
presupuesto fijo de tiempo de viaje sobre la red de transporte existente. La accesibilidad
---la facilidad de alcanzar destinos valiosos--- tiene una larga literatura conceptual y
metodológica \parencite{handy1997, geurs2004}; las medidas de oportunidades acumuladas
están entre sus formas más usadas por ser transparentes y requerir solo una red y una
distribución de oportunidades, a costa de un corte abrupto en el límite del presupuesto. El
presupuesto de tiempo combina aquí el tiempo de caminata de acceso, la espera esperada
derivada de las frecuencias y el tiempo a bordo (\cref{sec:w3}).

La contribución diagnóstica de la tesis vive en la intersección de demanda y
accesibilidad. En lugar de etiquetar zonas según si hoy tienen servicio ---el objetivo
circular criticado en \cref{sec:intro-circularidad}---, el marco define una \emph{brecha de
cobertura} como la demanda modelada frente a la accesibilidad medida, de modo que una
unidad puntúa alto solo cuando su demanda latente está genuinamente desatendida por la
oferta actual. Esto conecta con dos vertientes afines: la literatura de ``desiertos de
transporte'', que contrapone la oferta de transporte a la demanda a resolución espacial
fina \parencite{jiao2013deserts}, y la tradición de brecha por necesidad social, que
pondera la subprovisión por las necesidades sociales de la población afectada
\parencite{currie2010gaps}. Más cercano aún es el índice de necesidad de transporte basado
en accesibilidad de \textcite{almamun2011servicegaps}, que compara un índice de necesidad
contra un índice de accesibilidad a una resolución igualmente fina específicamente para
identificar brechas de servicio ---la misma estructura de comparación de dos índices que
implementa aquí el cociente de brecha de cobertura de esta tesis (\cref{eq:gap}). La
contribución del marco respecto a esta tradición de índice de brecha ya existente no es la
estructura de comparación en sí, que está bien establecida, sino propagar la brecha
resultante directamente como objetivo de aprendizaje supervisado (\cref{sec:marco-aprendizaje})
en lugar de detenerse en un mapa descriptivo. Dado que demanda y oferta se construyen de
forma independiente, esto produce una variable dependiente que un modelo supervisado puede
aprender sin memorizar simplemente la red existente.

\section{Análisis multicriterio y ponderación objetiva}
\label{sec:marco-multicriterio}

Convertir un conjunto de indicadores en una única puntuación de prioridad requiere pesos.
El análisis multicriterio está bien establecido en la evaluación de proyectos de transporte
\parencite{macharis2015}, pero los pesos elicitados de expertos o de encuestas importan
sesgo del analista y son difíciles de transferir entre ciudades. Esta tesis usa en cambio
ponderación objetiva, combinando el método CRITIC \parencite{diakoulaki1995critic} ---que
pondera un criterio por su intensidad de contraste y su independencia de los demás--- con
el método de entropía, que pondera por la dispersión de los valores de cada criterio medida
por la entropía de Shannon \parencite{shannon1948}. Ambos son basados en datos en el
sentido específico de que el paso final de asignación de pesos lee los pesos de la propia
matriz de indicadores, lo que hace la priorización reproducible y re-derivable desde cero
en una ciudad nueva sin volver a elicitar juicios de expertos \emph{para ese paso}. Esto no
elimina el juicio del analista de la metodología en su conjunto: qué catorce indicadores
entran a la matriz, cómo se agrupan en la taxonomía Nodo/Lugar/Personas, qué regla de
normalización (transformación logarítmica frente a min--max simple) recibe cada uno, y el
promedio de conjunto CRITIC--EWM con pesos iguales son todas decisiones del analista
tomadas río arriba, no leídas de los datos. El propio relato de
\cref{sec:datos-integridad} sobre una corrección de signo e imputación del índice de
marginación, que desplazó la puntuación media de equidad en varias décimas, es evidencia
directa de cuán consecuentes pueden ser estas decisiones río arriba, no ``objetivas'';
CRITIC/EWM retira el juicio experto del paso final de agregación, no de la metodología en
su conjunto. Esta calificación no es meramente una autocrítica: el análisis comparativo de
\textcite{mukhametzyanov2021critic} sobre la ponderación por entropía, CRITIC y desviación
estándar encuentra que la ponderación formal basada en la matriz de decisión no siempre se
comporta bien, lo cual confirma de forma independiente que el encuadre ``objetivo'' de
estos métodos necesita la misma calificación que esta tesis le aplica. Esta combinación de
ponderación de la familia CRITIC con un análisis independiente de brecha/accesibilidad no
es exclusiva de esta tesis: \textcite{rathod2024critic} aplican el mismo emparejamiento
---ponderación CRITIC-MCDM junto con un análisis de brecha de desempeño de accesibilidad de
transporte--- a Surat, India, un antecedente de 2024 para tratar la capa de priorización y
el diagnóstico de brecha como componentes relacionados pero distintos, tal como hace esta
tesis entre W3 y W4.

Un modo de falla específico y verificable que identifica la literatura anterior es la alta
correlación entre indicadores: la fórmula de peso de CRITIC descuenta a un criterio por
correlacionar con otros, pero el Método de Peso por Entropía no lo hace, de modo que un
promedio de conjunto de ambos solo corrige parcialmente por indicadores redundantes. Una
revisión de correlación de los catorce indicadores Nodo--Lugar--Personas (recalculada el
2026-09-09) encuentra tres pares por encima del umbral convencional de multicolinealidad
$|r|>0.8$---\code{p\_poi\_density\_n}~/~\code{p\_retail\_density\_n} ($r=0.93$),
\code{n\_intersections\_n}~/~\code{n\_street\_density\_n} ($r=0.828$), y
\code{pe\_population\_n}~/~\code{pe\_pop\_density\_n} ($r=0.812$)---sobre una correlación
media absoluta fuera de la diagonal de \num{0.335}. Ninguno de los tres pares es
sorprendente (cada uno empareja un conteo con una densidad o tasa estrechamente
relacionada), y la propia ponderación de CRITIC ya los descuenta; la verificación se
reporta aquí porque es el diagnóstico concreto que exige la literatura sobre modos de
falla, no porque revierta el resultado de la ponderación.

La equidad no se maneja manipulando estos pesos sino con un término aditivo explícito y
separable, construido a partir de índices de marginación y rezago social
\parencite{conapo_marginacion2020, coneval_rezago2020}, de modo que su contribución a la
puntuación final sea transparente y pueda reportarse con y sin él (\cref{eq:equity}). Este
tratamiento separable refleja una visión creciente de que la equidad en el transporte debe
ser un objetivo explícito y no un subproducto implícito
\parencite{martens2016, pereira2017justice}.

\section{Aprendizaje automático para ubicación de estaciones y corredores}
\label{sec:marco-aprendizaje}

Una literatura considerable aplica aprendizaje supervisado a la ubicación de transporte.
\textcite{niu2023rf} usan bosques aleatorios para predecir la idoneidad de estaciones, y
la línea de aprendizaje automático Nodo--Lugar revisada arriba \parencite{su2022nodeplace,
aminipishro2022nprt} combina modelos basados en árboles con valores
\textsc{shap} \parencite{lundberg2017shap} para interpretabilidad. El marco aquí usa las
mismas herramientas ---bosques aleatorios \parencite{breiman2001rf} y árboles potenciados
por gradiente \parencite{chen2016xgboost, ke2017lightgbm}, leídos a través de
\textsc{shap}--- pero se aparta de gran parte de esa literatura en la variable
dependiente. Cuando un modelo se entrena para predecir una etiqueta de ``tiene parada'', o
una clase de idoneidad que es a su vez función de la oferta actual, solo puede aprender a
reproducir el statu quo y reportará alta exactitud por una razón circular. La tesis hace
explícito este modo de falla: una versión anterior de su propia metodología se entrenó
sobre una etiqueta derivada de un agrupamiento y reportó una exactitud de recuperación casi
perfecta que no portaba información de planeación, y fue retirada en favor de un modelo
entrenado sobre el objetivo independiente de brecha de cobertura (\cref{sec:w3}).

Esta circularidad específica de la ubicación de transporte no es una observación aislada
sobre la historia de una sola metodología; es un caso particular de un mecanismo al que se
ha dado tratamiento formal preciso en tres literaturas más maduras, introducidas en
\cref{sec:intro-circularidad} y desarrolladas aquí. \textcite{kaufman2012leakage} nombran y
formalizan el modo de falla general en minería de datos ---la \emph{fuga}, el uso de
información derivada del objetivo a la que el modelo no debería tener acceso legítimo---
del cual una etiqueta de ubicación derivada de la oferta es una instancia.
\textcite{lakkaraju2017selectivelabels} con su \emph{problema de las etiquetas selectivas}
da el relato estadístico más agudo de por qué esta instancia particular de fuga es
inevitable por construcción y no meramente un defecto de calidad de datos: siempre que un
resultado se observa solo condicionado a una decisión previa (en su contexto, la libertad
bajo fianza; aquí, una decisión previa de ubicación), ninguna cantidad de datos o
sofisticación adicional de modelado recupera el resultado contrafactual para las unidades
donde la decisión fue en sentido contrario. Empíricamente, \textcite{lum2016predict} y, de
forma más rigurosa, \textcite{ensign2018runaway} muestran que los modelos de vigilancia
predictiva entrenados sobre dónde se envió previamente a la policía reproducen y pueden
amplificar ese patrón histórico en lugar de detectar dónde se concentra genuinamente el
crimen, un bucle de retroalimentación descontrolado con la estructura lógica idéntica
---entrenar sobre dónde ya ocurrió la intervención, y el modelo solo puede reaprender la
huella de la propia intervención--- y \textcite{flynn2022spacetime} plantean la misma
preocupación específicamente para la asignación espacial de infraestructura. Fundamentar el
argumento de circularidad de la tesis en este cuerpo de trabajo establecido y
transdisciplinario, en lugar de tratarlo como una observación ad hoc sobre la historia de
una metodología, es una de las correcciones hechas durante una revisión de literatura del
2026-09-09. El punto metodológico ---que las herramientas de interpretabilidad no pueden
rescatar un objetivo con etiquetado selectivo, y que el objetivo mismo debe construirse con
independencia de la decisión previa que de otro modo codificaría--- es una lección central
de este trabajo, y una con antecedentes mucho más allá de la planeación de transporte.

\section{Diseño de redes y generación de corredores}
\label{sec:marco-redes}

Generar corredores candidatos a partir de un conjunto de ubicaciones prioritarias es un
problema de diseño de redes, parte de la amplia literatura de diseño de redes de transporte
\parencite{guihaire2008, farahani2013review}, de la cual la revisión de
\textcite{kepaptsoglou2009review} ---que organiza el campo por objetivos de diseño,
parámetros del entorno de operación y método de solución--- sigue siendo el punto de
referencia estándar para clasificar un método nuevo como la combinación W5/W6 de esta
tesis. La tradición heurística relevante es el árbol
de Steiner en grafos ---conectar un conjunto de nodos requeridos (terminales) a costo
mínimo, posiblemente a través de nodos intermedios adicionales---, para el que
\textcite{takahashi1980} dan una aproximación clásica basada en caminos más cortos y
\textcite{kmb1981steiner} una rápida basada en árbol de expansión mínima, la familia de
parientes usada aquí. El diseño de transporte basado en árbol de expansión no es
puramente histórico: un estudio de caso de 2026 aplica la misma idea subyacente ---un
árbol de expansión refinado por búsqueda tabú que minimiza el total de
pasajeros-kilómetro--- a la red de autobuses de Canberra \parencite{suguira2026spanningtree},
evidencia de que la familia heurística sigue en uso activo para el (re)diseño de redes
reales, aunque al ser un preprint ese resultado debe ponderarse en consecuencia. Estas
heurísticas resultan atractivas para la generación de corredores
porque operan directamente sobre el grafo vial y producen geometría conectada y construible
a partir de anclas de demanda dispersas.

Su limitación, que la tesis caracteriza en detalle, es que un árbol construido para
minimizar el costo de conexión entre anclas no es el mismo objeto que un buen
\emph{corredor}. Un árbol de expansión convertido en una sola línea puede insertar saltos
fantasma entre ramas no adyacentes y ---más grave--- un filtro de factibilidad basado en el
rodeo entre extremos seleccionará pocas anclas casi colineales sin importar la demanda que
un corredor realmente atiende, confundiendo factibilidad con escasez de anclas.
\Cref{sec:w6} rearquitectura el generador en torno a un modelador de tronco por diámetro y
una métrica de factibilidad por rectitud respecto a las anclas, que miden lo correcto, y
\cref{ch:discusion} reporta tanto lo que esto corrige como la limitación residual que
persiste. Un punto de posicionamiento aparte, relativo a la taxonomía de
\textcite{kepaptsoglou2009review}, resulta más favorable al diseño de esta tesis de lo que
sugería una versión anterior de esta sección. La investigación académica de TNDP
multiobjetivo se ha movido hacia la exploración de Pareto multiobjetivo que evita imponer
pesos fijos del todo, y \textcite{silva2022mo} lo señalan directamente: es genuinamente
difícil para los expertos en transporte actuar sobre un frente de Pareto crudo, o combinar
objetivos de manera que se adapte con flexibilidad conforme cambian las prioridades, que es
precisamente por qué la práctica aplicada aún se apoya en la escalarización por suma
ponderada pese a su incapacidad de alcanzar regiones no convexas del espacio de
compensación. En lugar de tratar un enfoque como simplemente superado por el otro,
\textcite{silva2022mo} argumentan que los encuadres de uno y de varios objetivos se
combinan mejor de forma sinérgica. La puntuación compuesta en \cref{eq:objectives} hace
exactamente eso: combina $f_1$, eficiencia y $f_3$ con pesos fijos (0.50/0.25/0.25)
\emph{junto a}, no en lugar de, la clasificación de Pareto genuina (\cref{fig:pareto}), lo
cual da al lector tanto un número único accionable como el conjunto no dominado completo
para inspeccionar. Este híbrido es consistente con, y no está rezagado respecto a, cómo el
campo concilia actualmente la práctica académica y la aplicada en esta cuestión ---aunque,
como muestran \cref{sec:res-w6,sec:disc-generative}, tener ambos solo es útil si la
puntuación compuesta y la clasificación de Pareto realmente concuerdan, cosa que para el
conjunto de candidatos propio de esta tesis no ocurre.

\section{Vacío de investigación}
\label{sec:marco-vacio}

Las cuatro literaturas anteriores son individualmente maduras pero rara vez se combinan en
una única metodología reproducible que sea además \emph{no circular} y \emph{transferible}.
Los modelos de demanda estiman flujos pero pocas veces se propagan hacia un diagnóstico
explícito de brecha; los estudios Nodo--Lugar describen estaciones pero no generan
corredores; el trabajo de ubicación con aprendizaje automático frecuentemente entrena sobre
objetivos derivados de la oferta; y las heurísticas de generación de corredores rara vez se
evalúan contra necesidad genuina en lugar de contra la red que ya existe. Esto no quiere
decir que la mitad diagnóstica de esta afirmación carezca de antecedentes:
\textcite{foth2013equitable} construyen esencialmente la misma superposición de
accesibilidad frente a necesidad social usada aquí como diagnóstico espacial de equidad
para Toronto, y \textcite{camporeale2019equity} abordan directamente soluciones de diseño
de redes evaluadas contra principios de equidad horizontal y vertical en lugar de solo
costo o cobertura ---el mismo territorio que ocupa la capa generativa de esta tesis
(\cref{sec:w5,sec:w6}). Lo que no está combinado en ninguno de esos trabajos, ni en el
resto de la literatura revisada, es la cadena específica que cierra esta tesis: una
estimación de demanda basada solo en Nivel~1, alimentada a un reentrenamiento supervisado
no circular, acoplada a un generador de tronco por diámetro de árbol de expansión mínima
con puerta de factibilidad por rectitud respecto a las anclas, y toda la cadena
transferida de extremo a extremo entre zonas metropolitanas de tamaño contrastante. El
vacío que aborda esta tesis es, por tanto, más acotado que ``una metodología
integradora'' en sí misma ---ya existen diagnósticos integradores y generadores
sensibles a la equidad estrechamente relacionados--- y se establece mejor como: un marco
basado en la demanda
que (i) estima la demanda solo con datos de Nivel~1, (ii) diagnostica la brecha de
cobertura con un objetivo construido con independencia de la oferta y lo valida fuera de
muestra, (iii) genera y evalúa corredores candidatos por necesidad, no redundancia y
eficiencia de demanda, y (iv) se transfiere de extremo a extremo a zonas metropolitanas de
tamaño contrastante sin rediseño a la medida. La metodología del \cref{ch:metodologia} se
organiza para cerrar exactamente esta combinación más acotada, pero aún no abordada.
